\section{Số liệu đánh giá}
\subsection{Các chỉ số đánh giá (Evaluation Metrics)}

Phần này trình bày các chỉ số đánh giá được sử dụng trong toàn bộ hệ thống, bao phủ từ mức tín hiệu (speech processing) đến mức suy luận ngữ nghĩa (RAG). Các chỉ số được thiết kế nhằm phản ánh đầy đủ các tiêu chí: độ chính xác, khả năng bảo toàn thông tin, và hiệu quả truy xuất.

\subsubsection{Diarization Metrics}

\textbf{Vai trò:}
Đánh giá chất lượng của mô-đun định danh người nói, là bước nền tảng ảnh hưởng trực tiếp đến toàn bộ pipeline phía sau.

\textbf{Diarization Error Rate (DER):}

DER là chỉ số tiêu chuẩn để đo lường mức độ sai lệch giữa kết quả diarization và ground truth theo thời gian.

\[
DER = \frac{Miss + FalseAlarm + Confusion}{Total\ Speech\ Time}
\]

Trong đó:
\begin{itemize}
    \item \textbf{Miss (Missed Speech):} phần speech không được phát hiện
    \item \textbf{False Alarm:} phần không phải speech nhưng bị phát hiện nhầm
    \item \textbf{Confusion:} phần speech được gán sai người nói
\end{itemize}

Trong phạm vi hệ thống, phân tích tập trung chủ yếu vào:
\begin{itemize}
    \item \textbf{Miss:} ảnh hưởng trực tiếp đến việc mất thông tin đầu vào cho ASR
    \item \textbf{Confusion:} ảnh hưởng đến tính đúng đắn của metadata speaker, từ đó tác động đến retrieval và QA
\end{itemize}

\textbf{Speaker Matching (Label Alignment):}

Do nhãn speaker trong hệ thống mang tính ký hiệu (arbitrary), cần thực hiện bước ánh xạ nhãn giữa kết quả dự đoán và ground truth trước khi tính DER.

Bài toán này được mô hình hóa như một bài toán gán tối ưu và được giải bằng thuật toán Hungarian, dựa trên ma trận nhầm lẫn (confusion matrix) ở mức frame.

\textbf{Phân tích lỗi:}
\begin{itemize}
    \item Label oscillation (dao động nhãn)
    \item Over-merging (gộp nhiều speaker)
    \item Over-splitting (tách một speaker thành nhiều nhãn)
\end{itemize}

Các phân tích này được sử dụng như công cụ chẩn đoán nhằm cải thiện pipeline, không phải là chỉ số định lượng độc lập.

\textbf{Speaker-Aware Coverage (SAC):}

SAC là chỉ số bổ sung nhằm đo lường khả năng bảo toàn nội dung gắn với đúng speaker.

\begin{itemize}
    \item \textbf{Coverage:} tỷ lệ các từ/cụm từ có nghĩa trong Ground Truth xuất hiện trong output
    \item \textbf{Speaker Correctness:} tỷ lệ các từ được gán đúng speaker
\end{itemize}

SAC giúp đánh giá tác động của lỗi diarization lên các bước downstream, đặc biệt là truy xuất và tổng hợp thông tin.

\subsubsection{Summarization Metrics}

\textbf{Vai trò:}
Đánh giá khả năng của hệ thống trong việc rút gọn dữ liệu hội thoại dài thành biểu diễn cô đọng, đồng thời vẫn giữ được thông tin quan trọng phục vụ truy xuất.

\begin{itemize}
    \item \textbf{Coverage:} 
    Tỷ lệ các điểm thông tin quan trọng được giữ lại trong bản tóm tắt. 
    Ngưỡng yêu cầu tối thiểu là 70\%.

    \item \textbf{Compression Ratio:} 
    Tỷ lệ giữa độ dài văn bản tóm tắt và văn bản gốc, được kiểm soát trong khoảng 20\%–40\%.

    \item \textbf{Faithfulness:} 
    Đảm bảo nội dung tóm tắt không chứa thông tin không xuất hiện trong ngữ cảnh đầu vào (không hallucination).
\end{itemize}

\subsubsection{Retrieval Metrics}

\textbf{Vai trò:}
Đánh giá khả năng của hệ thống trong việc truy xuất đúng và ưu tiên các đoạn hội thoại chứa thông tin quan trọng.

\begin{itemize}
    \item \textbf{Hit Rate@K:} 
    Tỷ lệ truy vấn mà đoạn chứa đáp án đúng xuất hiện trong Top-$K$ kết quả.

    \item \textbf{Mean Reciprocal Rank (MRR):}
    \[
    MRR = \frac{1}{Q} \sum_{i=1}^{Q} \frac{1}{rank_i}
    \]
    Đánh giá mức độ ưu tiên của các kết quả đúng trong danh sách truy xuất.

    \item \textbf{Context Relevance:}
    Đo lường mức độ liên quan của ngữ cảnh truy xuất so với truy vấn, phản ánh khả năng loại bỏ thông tin nhiễu trước khi đưa vào mô hình sinh.
\end{itemize}

\subsubsection{Metadata Impact}

\textbf{Vai trò:}
Đánh giá mức độ đóng góp của metadata (speaker, timestamp) đối với hiệu năng truy xuất và hỏi đáp.

\begin{itemize}
    \item So sánh HitRate@K và MRR giữa hai cấu hình:
    \begin{itemize}
        \item Có sử dụng metadata
        \item Không sử dụng metadata
    \end{itemize}
\end{itemize}

Kết quả so sánh cho phép xác định vai trò của thông tin speaker trong việc giảm nhiễu ngữ nghĩa và cải thiện độ chính xác của hệ thống.